\setcounter{chapter}{10}
\chapter{Fubiniと変数変換}

\paragraph{本章の目標}
本章では，多変数のLebesgue積分を実際に計算するための2つの道具を導入する．
1つは積分順序を交換するFubiniの定理であり，
もう1つは座標を取り替える変数変換である．
特に
\begin{enumerate}
    \item 非負関数に対するTonelliの定理
    \item 可積分関数に対するFubiniの定理
    \item 平行移動・拡大縮小・せん断を用いた線形変数変換
    \item affine 変換
\end{enumerate}
を証明する．

\section{動機づけ}

2変数関数 $f(x,y)$ を積分するとき，
\[
\iint f(x,y)\,dx\,dy
\]
を
\[
\int\left(\int f(x,y)\,dy\right)dx
\]
あるいは
\[
\int\left(\int f(x,y)\,dx\right)dy
\]
として計算したくなる．
しかし，積分順序の交換はいつでも正しいわけではない．
その正当化を与えるのがTonelliの定理とFubiniの定理である．

また，積分を計算する際には
\[
x=au+b,\qquad
\begin{pmatrix}x\\ y\end{pmatrix}
=
T\begin{pmatrix}u\\ v\end{pmatrix}
\]
のように座標を取り替えることも多い．
変数変換公式は，
「面積や体積がどう伸び縮みするか」
を積分の中に正しく反映させる定理である．

\section{積測度と切片}

以後，$\RR^m$ 上のLebesgue測度を $\lambda_m$，
$\RR^n$ 上のLebesgue測度を $\lambda_n$ と書く．

\begin{definition}[積測度]\label{def:product-measure}
$\RR^m\times\RR^n$ 上の積 $\sigma$-加法族
\[
\calL_m\otimes \calL_n
\]
に対し，Carathéodory--Hahnの拡張定理
%（\cref{thm:caratheodory-hahn-extension}）
によって得られる測度
\[
\lambda_m\otimes\lambda_n
\]
で，
任意の可測矩形 $A\times B$ に対して
\[
(\lambda_m\otimes\lambda_n)(A\times B)=\lambda_m(A)\lambda_n(B)
\]
を満たすものを積測度という．
\end{definition}

\begin{remark}
本章では
\[
(\RR^m\times\RR^n,\calL_m\otimes\calL_n,\lambda_m\otimes\lambda_n)
\]
の上で積分を考える．
直感的には
\[
d\lambda_m(x)\,d\lambda_n(y)
\]
と書いてよい．
\end{remark}

\begin{definition}[切片]\label{def:sections}
$E\subset \RR^m\times\RR^n$ と
$f:\RR^m\times\RR^n\to\eRR$ に対して，
各 $x\in\RR^m$，$y\in\RR^n$ について
\[
E_x:=\{y\in\RR^n\mid (x,y)\in E\},
\qquad
E^y:=\{x\in\RR^m\mid (x,y)\in E\}
\]
を集合 $E$ の $x$-切片，$y$-切片という．
また
\[
f_x(y):=f(x,y),
\qquad
f^y(x):=f(x,y)
\]
を関数 $f$ の切片という．
\end{definition}

\begin{proposition}[切片の可測性]\label{prop:measurable-sections}
$E\in \calL_m\otimes\calL_n$ とする．
このとき任意の $x\in\RR^m$ と任意の $y\in\RR^n$ に対して
\[
E_x\in \calL_n,
\qquad
E^y\in \calL_m
\]
が成り立つ．
\end{proposition}
\begin{proof}
$x\in\RR^m$ を固定し，
\[
\mathcal C_x:=\{E\in \calL_m\otimes\calL_n \mid E_x\in\calL_n\}
\]
とおく．
まず $\mathcal C_x$ は $\sigma$-加法族である．
実際，
\[
\emptyset_x=\emptyset
\]
だから $\emptyset\in\mathcal C_x$ であり，
\[
(\RR^m\times\RR^n\setminus E)_x=\RR^n\setminus E_x
\]
だから補集合に関して閉じている．
また
\[
\left(\bigcup_{k=1}^{\infty}E_k\right)_x=\bigcup_{k=1}^{\infty}(E_k)_x
\]
だから可算和に関しても閉じている．

次に，任意の可測矩形 $A\times B$ に対して
\[
(A\times B)_x
=
\begin{cases}
B & (x\in A),\\
\emptyset & (x\notin A)
\end{cases}
\]
であるから $A\times B\in\mathcal C_x$ である．
したがって，矩形が生成する $\sigma$-加法族全体
\[
\calL_m\otimes\calL_n
\]
が $\mathcal C_x$ に含まれる．
ゆえに $E_x$ は可測である．

$E^y$ についても同様である．
\end{proof}

\begin{corollary}[切片関数の可測性]\label{cor:measurable-function-sections}
$f:\RR^m\times\RR^n\to\eRR$ を
\[
\calL_m\otimes\calL_n
\]
に関して可測とする．
このとき任意の $x\in\RR^m$ と任意の $y\in\RR^n$ に対して
\[
f_x:\RR^n\to\eRR,
\qquad
f^y:\RR^m\to\eRR
\]
は可測である．
\end{corollary}
\begin{proof}
任意の $a\in\RR$ に対して
\[
\{y\in\RR^n\mid f_x(y)>a\}
=
\{(u,v)\in\RR^m\times\RR^n\mid f(u,v)>a\}_x
\]
である．
右辺に現れる集合は $f$ の可測性より
\[
\calL_m\otimes\calL_n
\]
に属するから，\cref{prop:measurable-sections}よりその $x$-切片は可測である．
したがって\cref{thm:lebesgue-measurable-function-criterion}より $f_x$ は可測である．

$f^y$ についても同様である．
\end{proof}

\section{集合に対するFubini型公式}

\begin{theorem}[有界矩形の中の集合]\label{thm:set-fubini-bounded-rectangle}
$I\subset\RR^m$，$J\subset\RR^n$ を有界閉区間とする．
このとき，任意の
\[
E\in \calL_m(I)\otimes\calL_n(J)
\]
に対して
\[
x\longmapsto \lambda_n(E_x)
\]
は $I$ 上可測であり，
\[
(\lambda_m\otimes\lambda_n)(E)=\int_I \lambda_n(E_x)\,d\lambda_m(x)
\]
が成り立つ．
ただし
\[
\calL_m(I):=\{A\cap I\mid A\in\calL_m\},
\qquad
\calL_n(J):=\{B\cap J\mid B\in\calL_n\}
\]
とする．
\end{theorem}
\begin{proof}
\[
\mathcal C
:=
\left\{
E\in \calL_m(I)\otimes\calL_n(J)
\ \middle|\
x\mapsto \lambda_n(E_x)\text{ は }I\text{ 上可測で， }
(\lambda_m\otimes\lambda_n)(E)=\int_I \lambda_n(E_x)\,d\lambda_m(x)
\right\}
\]
とおく．
まず $\mathcal C$ が $I\times J$ 上の $\sigma$-加法族であることを示す．

\textbf{(1) 空集合}\quad
\[
\emptyset_x=\emptyset
\]
だから
\[
\lambda_n(\emptyset_x)=0
\]
であり，
\[
(\lambda_m\otimes\lambda_n)(\emptyset)=0=\int_I 0\,d\lambda_m(x)
\]
である．
したがって $\emptyset\in\mathcal C$ である．

\textbf{(2) 補集合}\quad
$E\in\mathcal C$ とする．
各 $x\in I$ に対して
\[
\bigl((I\times J)\setminus E\bigr)_x=J\setminus E_x
\]
だから
\[
\lambda_n\bigl(((I\times J)\setminus E)_x\bigr)=\lambda_n(J)-\lambda_n(E_x)
\]
である．
右辺は可測関数だから，左辺も可測である．
さらに
\begin{align*}
\int_I \lambda_n\bigl(((I\times J)\setminus E)_x\bigr)\,d\lambda_m(x)
&=
\int_I \bigl(\lambda_n(J)-\lambda_n(E_x)\bigr)\,d\lambda_m(x) \\
&=
\lambda_m(I)\lambda_n(J)-\int_I \lambda_n(E_x)\,d\lambda_m(x) \\
&=
(\lambda_m\otimes\lambda_n)(I\times J)-(\lambda_m\otimes\lambda_n)(E) \\
&=
(\lambda_m\otimes\lambda_n)\bigl((I\times J)\setminus E\bigr).
\end{align*}
したがって $(I\times J)\setminus E\in\mathcal C$ である．

\textbf{(3) 互いに素な可算和}\quad
$E_k\in\mathcal C$ が互いに素であるとする．
すると各 $x\in I$ に対して
\[
\left(\bigcup_{k=1}^{\infty}E_k\right)_x=\bigcup_{k=1}^{\infty}(E_k)_x
\]
であり，右辺は互いに素である．
よって
\[
\lambda_n\left(\left(\bigcup_{k=1}^{\infty}E_k\right)_x\right)
=
\sum_{k=1}^{\infty}\lambda_n((E_k)_x).
\]
右辺は可測関数列の和だから可測である．
また\cref{thm:monotone-convergence}より
\begin{align*}
\int_I \lambda_n\left(\left(\bigcup_{k=1}^{\infty}E_k\right)_x\right)\,d\lambda_m(x)
&=
\int_I \sum_{k=1}^{\infty}\lambda_n((E_k)_x)\,d\lambda_m(x) \\
&=
\sum_{k=1}^{\infty}\int_I \lambda_n((E_k)_x)\,d\lambda_m(x) \\
&=
\sum_{k=1}^{\infty}(\lambda_m\otimes\lambda_n)(E_k) \\
&=
(\lambda_m\otimes\lambda_n)\left(\bigcup_{k=1}^{\infty}E_k\right).
\end{align*}
したがって $\bigcup_k E_k\in\mathcal C$ である．

以上より $\mathcal C$ は $\sigma$-加法族である．

次に矩形 $A\times B$ で
$A\in\calL_m(I)$，$B\in\calL_n(J)$ とする．
このとき
\[
(A\times B)_x
=
\begin{cases}
B & (x\in A),\\
\emptyset & (x\notin A)
\end{cases}
\]
だから
\[
\lambda_n((A\times B)_x)=\lambda_n(B)1_A(x)
\]
である．
したがって
\[
\int_I \lambda_n((A\times B)_x)\,d\lambda_m(x)
=
\lambda_n(B)\lambda_m(A)
=
(\lambda_m\otimes\lambda_n)(A\times B).
\]
よって可測矩形はすべて $\mathcal C$ に属する．

可測矩形は
\[
\calL_m(I)\otimes\calL_n(J)
\]
を生成するから，
\[
\mathcal C=\calL_m(I)\otimes\calL_n(J)
\]
である．
\end{proof}

\begin{theorem}[一般の集合]\label{thm:set-fubini-general}
任意の
\[
E\in\calL_m\otimes\calL_n
\]
に対して
\[
x\longmapsto \lambda_n(E_x)
\]
は $\RR^m$ 上可測であり，
\[
(\lambda_m\otimes\lambda_n)(E)=\int_{\RR^m}\lambda_n(E_x)\,d\lambda_m(x)
\]
が成り立つ．
同様に
\[
y\longmapsto \lambda_m(E^y)
\]
は $\RR^n$ 上可測であり，
\[
(\lambda_m\otimes\lambda_n)(E)=\int_{\RR^n}\lambda_m(E^y)\,d\lambda_n(y)
\]
が成り立つ．
\end{theorem}
\begin{proof}
まず $x$ に関する式を示す．
各 $k\in\NN$ に対して
\[
I_k:=[-k,k]^m,\qquad J_k:=[-k,k]^n,
\qquad
E_k:=E\cap (I_k\times J_k)
\]
とおく．
すると
\[
E_1\subset E_2\subset \cdots,\qquad
\bigcup_{k=1}^{\infty}E_k=E
\]
である．
また各 $k$ に対して\cref{thm:set-fubini-bounded-rectangle}より
\[
x\mapsto \lambda_n((E_k)_x)
\]
は可測で，
\[
(\lambda_m\otimes\lambda_n)(E_k)=\int_{\RR^m}\lambda_n((E_k)_x)\,d\lambda_m(x)
\]
が成り立つ．

固定した $x\in\RR^m$ に対して
\[
(E_k)_x=E_x\cap J_k
\]
だから
\[
\lambda_n((E_k)_x)\uparrow \lambda_n(E_x)
\]
である．
したがって\cref{thm:monotone-convergence}より
\begin{align*}
\int_{\RR^m}\lambda_n(E_x)\,d\lambda_m(x)
&=
\lim_{k\to\infty}\int_{\RR^m}\lambda_n((E_k)_x)\,d\lambda_m(x) \\
&=
\lim_{k\to\infty}(\lambda_m\otimes\lambda_n)(E_k) \\
&=
(\lambda_m\otimes\lambda_n)(E).
\end{align*}
可測性も極限として従う．

$y$ に関する式も同様に示される．
\end{proof}

\section{Tonelliの定理}

\begin{theorem}[単関数に対するTonelli]\label{thm:tonelli-simple-functions}
非負単関数
\[
s:\RR^m\times\RR^n\to[0,\infty)
\]
に対して，
\[
x\longmapsto \int_{\RR^n}s(x,y)\,d\lambda_n(y)
\]
は可測であり，
\[
\int_{\RR^m\times\RR^n}s\,d(\lambda_m\otimes\lambda_n)
=
\int_{\RR^m}\left(\int_{\RR^n}s(x,y)\,d\lambda_n(y)\right)\,d\lambda_m(x)
\]
が成り立つ．
同様に
\[
\int_{\RR^m\times\RR^n}s\,d(\lambda_m\otimes\lambda_n)
=
\int_{\RR^n}\left(\int_{\RR^m}s(x,y)\,d\lambda_m(x)\right)\,d\lambda_n(y)
\]
も成り立つ．
\end{theorem}
\begin{proof}
$s$ は有限個の非負実数 $a_1,\dots,a_N$ と互いに素な可測集合
$E_1,\dots,E_N\in\calL_m\otimes\calL_n$ を用いて
\[
s=\sum_{k=1}^N a_k1_{E_k}
\]
と書ける．
すると各 $x\in\RR^m$ に対して
\[
s_x(y)=\sum_{k=1}^N a_k1_{(E_k)_x}(y)
\]
だから
\[
\int_{\RR^n}s(x,y)\,d\lambda_n(y)
=
\sum_{k=1}^N a_k \lambda_n((E_k)_x).
\]
右辺は可測関数の有限和だから可測である．

また集合に対するFubini型公式（\cref{thm:set-fubini-general}）より
\[
(\lambda_m\otimes\lambda_n)(E_k)
=
\int_{\RR^m}\lambda_n((E_k)_x)\,d\lambda_m(x)
\]
だから
\begin{align*}
\int_{\RR^m\times\RR^n}s\,d(\lambda_m\otimes\lambda_n)
&=
\sum_{k=1}^N a_k(\lambda_m\otimes\lambda_n)(E_k) \\
&=
\sum_{k=1}^N a_k\int_{\RR^m}\lambda_n((E_k)_x)\,d\lambda_m(x) \\
&=
\int_{\RR^m}\sum_{k=1}^N a_k\lambda_n((E_k)_x)\,d\lambda_m(x) \\
&=
\int_{\RR^m}\left(\int_{\RR^n}s(x,y)\,d\lambda_n(y)\right)\,d\lambda_m(x).
\end{align*}
$y$ についての式も同様である．
\end{proof}

\begin{theorem}[Tonelliの定理]\label{thm:tonelli}
$f:\RR^m\times\RR^n\to[0,\infty]$ を
\[
\calL_m\otimes\calL_n
\]
に関して可測な非負関数とする．
このとき
\[
x\longmapsto \int_{\RR^n}f(x,y)\,d\lambda_n(y)
\]
は可測であり，
\[
\int_{\RR^m\times\RR^n}f\,d(\lambda_m\otimes\lambda_n)
=
\int_{\RR^m}\left(\int_{\RR^n}f(x,y)\,d\lambda_n(y)\right)\,d\lambda_m(x)
\]
が成り立つ．
同様に
\[
\int_{\RR^m\times\RR^n}f\,d(\lambda_m\otimes\lambda_n)
=
\int_{\RR^n}\left(\int_{\RR^m}f(x,y)\,d\lambda_m(x)\right)\,d\lambda_n(y)
\]
も成り立つ．
\end{theorem}
\begin{proof}
\cref{lem:simple-function-approximation}により，
非負単関数列 $\{s_k\}$ で
\[
0\le s_1\le s_2\le\cdots,\qquad s_k(x,y)\uparrow f(x,y)
\]
を満たすものを取る．

各 $x$ に対して
\[
(s_k)_x(y)\uparrow f_x(y)
\]
だから，\cref{thm:monotone-convergence}より
\[
\int_{\RR^n}s_k(x,y)\,d\lambda_n(y)
\uparrow
\int_{\RR^n}f(x,y)\,d\lambda_n(y).
\]
左辺は\cref{thm:tonelli-simple-functions}により可測関数だから，極限も可測である．

さらに\cref{thm:tonelli-simple-functions}より
\[
\int_{\RR^m\times\RR^n}s_k\,d(\lambda_m\otimes\lambda_n)
=
\int_{\RR^m}\left(\int_{\RR^n}s_k(x,y)\,d\lambda_n(y)\right)\,d\lambda_m(x)
\]
である．
したがって再び\cref{thm:monotone-convergence}より
\begin{align*}
\int_{\RR^m\times\RR^n}f\,d(\lambda_m\otimes\lambda_n)
&=
\lim_{k\to\infty}\int_{\RR^m\times\RR^n}s_k\,d(\lambda_m\otimes\lambda_n) \\
&=
\lim_{k\to\infty}
\int_{\RR^m}\left(\int_{\RR^n}s_k(x,y)\,d\lambda_n(y)\right)\,d\lambda_m(x) \\
&=
\int_{\RR^m}
\left(
\lim_{k\to\infty}\int_{\RR^n}s_k(x,y)\,d\lambda_n(y)
\right)
d\lambda_m(x) \\
&=
\int_{\RR^m}\left(\int_{\RR^n}f(x,y)\,d\lambda_n(y)\right)\,d\lambda_m(x).
\end{align*}
$y$ についての式も同様である．
\end{proof}

\section{Fubiniの定理}

\begin{theorem}[Fubiniの定理]\label{thm:fubini}
$f:\RR^m\times\RR^n\to\eRR$ を
\[
\calL_m\otimes\calL_n
\]
に関して可測で，
\[
\int_{\RR^m\times\RR^n}|f|\,d(\lambda_m\otimes\lambda_n)<\infty
\]
を満たす可積分関数とする．
このとき次が成り立つ．
\begin{enumerate}
    \item ほとんどすべての $x\in\RR^m$ に対して，$f_x$ は $\RR^n$ 上可積分である
    \item
    \[
    g(x):=\int_{\RR^n}f(x,y)\,d\lambda_n(y)
    \]
    は $\lambda_m$-a.e.\ 定義される可積分関数である
    \item
    \[
    \int_{\RR^m\times\RR^n}f\,d(\lambda_m\otimes\lambda_n)
    =
    \int_{\RR^m}g(x)\,d\lambda_m(x)
    \]
    が成り立つ
\end{enumerate}
同様に，ほとんどすべての $y\in\RR^n$ に対して $f^y$ は可積分であり，
\[
\int_{\RR^m\times\RR^n}f\,d(\lambda_m\otimes\lambda_n)
=
\int_{\RR^n}\left(\int_{\RR^m}f(x,y)\,d\lambda_m(x)\right)\,d\lambda_n(y)
\]
が成り立つ．
\end{theorem}
\begin{proof}
\cref{thm:tonelli}を $|f|$ に適用すると
\[
h(x):=\int_{\RR^n}|f(x,y)|\,d\lambda_n(y)
\]
は可測で，
\[
\int_{\RR^m}h(x)\,d\lambda_m(x)
=
\int_{\RR^m\times\RR^n}|f|\,d(\lambda_m\otimes\lambda_n)
<
\infty
\]
である．
したがって
\[
h(x)<\infty
\]
がほとんどすべての $x$ で成り立つ．
このときそのような $x$ では $f_x$ は可積分である．
これで (1) が従う．

次に
\[
g_+(x):=\int_{\RR^n}f^+(x,y)\,d\lambda_n(y),
\qquad
g_-(x):=\int_{\RR^n}f^-(x,y)\,d\lambda_n(y)
\]
とおく．
\cref{thm:tonelli}より $g_+,g_-$ は可測であり，
\[
\int_{\RR^m}g_+(x)\,d\lambda_m(x)=\int f^+\,d(\lambda_m\otimes\lambda_n),
\qquad
\int_{\RR^m}g_-(x)\,d\lambda_m(x)=\int f^-\,d(\lambda_m\otimes\lambda_n)
\]
が成り立つ．
しかも
\[
g_+(x)+g_-(x)=h(x)<\infty
\]
がほとんどすべての $x$ で成り立つから，
\[
g(x):=g_+(x)-g_-(x)
\]
は a.e.\ で意味をもち，可測である．
また
\[
|g(x)|\le h(x)
\]
が a.e.\ で成り立つので，$g$ は可積分である．
これで (2) が従う．

最後に
\begin{align*}
\int_{\RR^m}g(x)\,d\lambda_m(x)
&=
\int_{\RR^m}g_+(x)\,d\lambda_m(x)-\int_{\RR^m}g_-(x)\,d\lambda_m(x) \\
&=
\int_{\RR^m\times\RR^n}f^+\,d(\lambda_m\otimes\lambda_n)
-\int_{\RR^m\times\RR^n}f^-\,d(\lambda_m\otimes\lambda_n) \\
&=
\int_{\RR^m\times\RR^n}f\,d(\lambda_m\otimes\lambda_n).
\end{align*}
これで (3) が従う．

$y$ についての主張も全く同様である．
\end{proof}

\begin{corollary}[積分順序の交換]\label{cor:iterated-integral-exchange}
$f:\RR^m\times\RR^n\to\eRR$ が可測で，
\[
\int_{\RR^m}\left(\int_{\RR^n}|f(x,y)|\,d\lambda_n(y)\right)\,d\lambda_m(x)<\infty
\]
とする．
このとき $f$ は積測度に関して可積分であり，
\[
\int_{\RR^m}\left(\int_{\RR^n}f(x,y)\,d\lambda_n(y)\right)\,d\lambda_m(x)
=
\int_{\RR^n}\left(\int_{\RR^m}f(x,y)\,d\lambda_m(x)\right)\,d\lambda_n(y)
\]
が成り立つ．
\end{corollary}
\begin{proof}
左辺の仮定は\cref{thm:tonelli}により
\[
\int_{\RR^m\times\RR^n}|f|\,d(\lambda_m\otimes\lambda_n)<\infty
\]
と同値である．
したがって\cref{thm:fubini}を適用すればよい．
\end{proof}

\begin{example}[三角形の面積]
\[
E:=\{(x,y)\in[0,1]^2\mid x+y\le 1\}
\]
とする．
このとき
\[
(\lambda_1\otimes \lambda_1)(E)=\frac12
\]
である．
\end{example}
\begin{proof}
各 $x\in[0,1]$ に対して
\[
E_x=[0,1-x]
\]
である．
したがって
\[
\lambda_1(E_x)=1-x.
\]
よって集合に対する公式より
\[
(\lambda_1\otimes \lambda_1)(E)
=
\int_0^1(1-x)\,dx
=
\frac12.
\]
\end{proof}

\section{1変数の affine 変数変換}

\begin{proposition}[1次元の集合の affine 変換]\label{prop:one-dimensional-affine-set-transform}
$A\subset\RR$ をLebesgue可測集合とし，$a\in\RR\setminus\{0\}$，$b\in\RR$ とする．
このとき
\[
aA+b:=\{ax+b\mid x\in A\}
\]
はLebesgue可測であり，
\[
\lambda_1(aA+b)=|a|\,\lambda_1(A)
\]
が成り立つ．
\end{proposition}
\begin{proof}
まずLebesgue測度の平行移動不変性より
%（\cref{prop:lebesgue-measure-translation-invariant}）より
\[
\lambda_1(aA+b)=\lambda_1(aA)
\]
であるから，$b=0$ の場合を示せば十分である．

以後 $b=0$ とする．
Lebesgue外測度の定義から，任意の集合 $B\subset\RR$ に対して
\[
\lambda^*(aB)=|a|\,\lambda^*(B)
\]
が成り立つ．
実際，$B\subset\bigcup_n I_n$ なら
\[
aB\subset\bigcup_n aI_n
\]
であり，区間の長さは
\[
|aI_n|=|a|\,|I_n|
\]
だから
\[
\lambda^*(aB)\le |a|\sum_n |I_n|.
\]
被覆の取り方について下限を取れば
\[
\lambda^*(aB)\le |a|\lambda^*(B)
\]
を得る．
逆向きの不等式は $B=a^{-1}(aB)$ に同じ議論を適用すれば
\[
\lambda^*(B)\le \frac1{|a|}\lambda^*(aB)
\]
となることから従う．

次に $A$ が可測であることから，任意の $S\subset\RR$ に対して
\[
\lambda^*(S)=\lambda^*(S\cap A)+\lambda^*(S\setminus A)
\]
である．
ここで任意の $T\subset\RR$ に対して $S=a^{-1}T$ と置くと
\[
\lambda^*(a^{-1}T)
=
\lambda^*(a^{-1}T\cap A)+\lambda^*(a^{-1}T\setminus A).
\]
両辺に $|a|$ を掛けて外測度の拡大縮小公式を用いると
\[
\lambda^*(T)
=
\lambda^*(T\cap aA)+\lambda^*(T\setminus aA)
\]
を得る．
したがって $aA$ はLebesgue可測である．

最後に可測集合上では測度と外測度が一致するから
\[
\lambda_1(aA)=\lambda^*(aA)=|a|\,\lambda^*(A)=|a|\,\lambda_1(A)
\]
である．
\end{proof}

\begin{theorem}[1変数の affine 変数変換]\label{thm:one-dimensional-affine-change-variables}
$f:\RR\to[0,\infty]$ を非負可測関数とし，
$a\in\RR\setminus\{0\}$，$b\in\RR$ とする．
このとき
\[
\int_{\RR}f(au+b)\,d\lambda_1(u)
=
\frac{1}{|a|}\int_{\RR}f(x)\,d\lambda_1(x)
\]
が成り立つ．
\end{theorem}
\begin{proof}
まず指示関数
\[
f=1_A
\]
について示す．
このとき
\[
1_A(au+b)=1_{(A-b)/a}(u)
\]
だから
\[
\int_{\RR}1_A(au+b)\,d\lambda_1(u)
=
\lambda_1\left(\frac{A-b}{a}\right)
=
\frac{1}{|a|}\lambda_1(A)
=
\frac{1}{|a|}\int_{\RR}1_A(x)\,d\lambda_1(x).
\]

次に非負単関数
\[
s=\sum_{k=1}^N c_k1_{A_k}
\]
については，線形性より
\[
\int_{\RR}s(au+b)\,d\lambda_1(u)
=
\sum_{k=1}^N c_k\int_{\RR}1_{A_k}(au+b)\,d\lambda_1(u)
=
\frac{1}{|a|}\sum_{k=1}^N c_k\lambda_1(A_k)
=
\frac{1}{|a|}\int_{\RR}s(x)\,d\lambda_1(x).
\]

最後に一般の非負可測関数 $f$ を取る．
単関数列 $s_N\uparrow f$ を取れば
\[
s_N(au+b)\uparrow f(au+b).
\]
したがって\cref{thm:monotone-convergence}より
\begin{align*}
\int_{\RR}f(au+b)\,d\lambda_1(u)
&=
\lim_{N\to\infty}\int_{\RR}s_N(au+b)\,d\lambda_1(u) \\
&=
\lim_{N\to\infty}\frac{1}{|a|}\int_{\RR}s_N(x)\,d\lambda_1(x) \\
&=
\frac{1}{|a|}\int_{\RR}f(x)\,d\lambda_1(x).
\end{align*}
\end{proof}

\begin{corollary}\label{cor:one-dimensional-affine-change-variables-integrable}
$f:\RR\to\eRR$ が可積分ならば
\[
\int_{\RR}f(au+b)\,d\lambda_1(u)
=
\frac{1}{|a|}\int_{\RR}f(x)\,d\lambda_1(x)
\]
が成り立つ．
\end{corollary}
\begin{proof}
\cref{thm:one-dimensional-affine-change-variables}を $f^+$ と $f^-$ に適用して差を取ればよい．
\end{proof}

\section{基本的な線形変数変換}

\begin{lemma}[座標の入れ替え]\label{lem:change-variables-coordinate-permutation}
$f:\RR^d\to[0,\infty]$ を非負可測関数とし，
$P$ を座標の置換に対応する置換行列とする．
このとき
\[
\int_{\RR^d}f(Pu)\,d\lambda_d(u)=\int_{\RR^d}f(x)\,d\lambda_d(x)
\]
が成り立つ．
\end{lemma}
\begin{proof}
置換は隣り合う2座標の交換の合成で書けるから，
第1座標と第2座標を交換する場合を示せば十分である．

$u=(u_1,u_2,u')$ と書く．
\cref{thm:tonelli}より
\begin{align*}
\int_{\RR^d}f(u_2,u_1,u')\,d\lambda_d(u)
&=
\int_{\RR^{d-2}}\int_{\RR}\int_{\RR}f(u_2,u_1,u')\,d\lambda_1(u_1)\,d\lambda_1(u_2)\,d\lambda_{d-2}(u') \\
&=
\int_{\RR^{d-2}}\int_{\RR}\int_{\RR}f(v,u,u')\,d\lambda_1(u)\,d\lambda_1(v)\,d\lambda_{d-2}(u').
\end{align*}
ここで積分変数 $u,v$ はダミー変数だから，名前を入れ替えて
\[
\int_{\RR^{d-2}}\int_{\RR}\int_{\RR}f(u,v,u')\,d\lambda_1(v)\,d\lambda_1(u)\,d\lambda_{d-2}(u')
\]
となる．
これは再び\cref{thm:tonelli}より
\[
\int_{\RR^d}f(x)\,d\lambda_d(x)
\]
に等しい．
\end{proof}

\begin{lemma}[対角行列]\label{lem:change-variables-diagonal}
$f:\RR^d\to[0,\infty]$ を非負可測関数とし，
\[
D=\mathrm{diag}(a_1,\dots,a_d)
\]
を対角行列とする．
ここで $a_i\ne 0$ とする．
このとき
\[
\int_{\RR^d}f(Du)\,d\lambda_d(u)
=
\frac{1}{|a_1\cdots a_d|}\int_{\RR^d}f(x)\,d\lambda_d(x)
\]
が成り立つ．
\end{lemma}
\begin{proof}
\cref{thm:tonelli}により
\[
\int_{\RR^d}f(a_1u_1,\dots,a_du_d)\,d\lambda_d(u)
\]
を反復積分として計算できる．
まず $u_1$ について1変数の affine 変数変換を適用すると
\begin{align*}
\int_{\RR}f(a_1u_1,a_2u_2,\dots,a_du_d)\,d\lambda_1(u_1)
=
\frac1{|a_1|}
\int_{\RR}f(x_1,a_2u_2,\dots,a_du_d)\,d\lambda_1(x_1).
\end{align*}
これを $u_2,\dots,u_d$ に関する外側の積分に代入する．
同じ操作を $u_2,\dots,u_d$ について順に繰り返すと
\[
\int_{\RR^d}f(Du)\,d\lambda_d(u)
=
\frac{1}{|a_1\cdots a_d|}\int_{\RR^d}f(x)\,d\lambda_d(x)
\]
を得る．
\end{proof}

\begin{lemma}[せん断変換]\label{lem:change-variables-shear}
$f:\RR^d\to[0,\infty]$ を非負可測関数とし，
\[
S(u_1,\dots,u_d)=(u_1+cu_2,u_2,\dots,u_d)
\]
とする．
このとき
\[
\int_{\RR^d}f(Su)\,d\lambda_d(u)=\int_{\RR^d}f(x)\,d\lambda_d(x)
\]
が成り立つ．
\end{lemma}
\begin{proof}
\cref{thm:tonelli}より
\begin{align*}
\int_{\RR^d}f(u_1+cu_2,u_2,\dots,u_d)\,d\lambda_d(u)
=
\int_{\RR^{d-1}}
\left(
\int_{\RR}f(u_1+cu_2,u_2,\dots,u_d)\,d\lambda_1(u_1)
\right)
d\lambda_{d-1}(u_2,\dots,u_d).
\end{align*}
内側の積分に1変数の affine 変数変換（\cref{thm:one-dimensional-affine-change-variables}）を $a=1$, $b=cu_2$ として適用すると
\[
\int_{\RR}f(u_1+cu_2,u_2,\dots,u_d)\,d\lambda_1(u_1)
=
\int_{\RR}f(x_1,u_2,\dots,u_d)\,d\lambda_1(x_1)
\]
である．
したがって右辺は
\[
\int_{\RR^d}f(x)\,d\lambda_d(x)
\]
に等しい．
\end{proof}

\begin{proposition}[可逆行列の分解]\label{prop:invertible-matrix-elementary-decomposition}
任意の可逆行列 $T\in GL(d,\RR)$ は，
有限個の
\begin{enumerate}
    \item 座標の置換
    \item 対角行列
    \item せん断変換
\end{enumerate}
に対応する行列の積として表せる．
\end{proposition}
\begin{proof}
Gaussian 消去法を $T$ に施すと，
有限回の基本行列操作によって単位行列 $I$ に変形できる．
すなわち，ある基本行列 $E_1,\dots,E_N$ が存在して
\[
E_N\cdots E_1T=I
\]
となる．
したがって
\[
T=E_1^{-1}\cdots E_N^{-1}.
\]
ここで基本行列の逆行列は，再び同じ3種類のいずれかである．
実際，
\begin{enumerate}
    \item 行の交換の逆は同じ交換
    \item 行の非零倍の逆は逆数倍
    \item 1行に他の行の定数倍を加える操作の逆は，その符号を反転した操作
\end{enumerate}
である．
よって主張が従う．
\end{proof}

\section{一般の線形変数変換}

\begin{theorem}[線形変数変換]\label{thm:linear-change-of-variables}
$T\in GL(d,\RR)$ とし，
$f:\RR^d\to[0,\infty]$ を非負可測関数とする．
このとき
\[
\int_{\RR^d}f(Tu)\,d\lambda_d(u)
=
\frac{1}{|\det T|}\int_{\RR^d}f(x)\,d\lambda_d(x)
\]
が成り立つ．
\end{theorem}
\begin{proof}
\cref{prop:invertible-matrix-elementary-decomposition}により
\[
T=E_1E_2\cdots E_N
\]
と書ける．
ここで各 $E_j$ は置換行列，対角行列，せん断変換のいずれかとする．

$N=1$ の場合は
\cref{lem:change-variables-coordinate-permutation,lem:change-variables-diagonal,lem:change-variables-shear}
で証明済みである．
一般の $N$ については帰納法で示す．
すでに
\[
R:=E_2\cdots E_N
\]
に対して
\[
\int_{\RR^d}f(E_1Ru)\,d\lambda_d(u)
\]
を考える．
関数
\[
g(v):=f(E_1v)
\]
は非負可測である．
帰納法の仮定を $g$ と $R$ に適用すると
\[
\int_{\RR^d}f(E_1Ru)\,d\lambda_d(u)
=
\frac{1}{|\det R|}\int_{\RR^d}f(E_1v)\,d\lambda_d(v).
\]
さらに $E_1$ に対する補題を適用すれば
\[
\int_{\RR^d}f(E_1v)\,d\lambda_d(v)
=
\frac{1}{|\det E_1|}\int_{\RR^d}f(x)\,d\lambda_d(x).
\]
両式を合わせると
\[
\int_{\RR^d}f(Tu)\,d\lambda_d(u)
=
\frac{1}{|\det E_1||\det R|}\int_{\RR^d}f(x)\,d\lambda_d(x)
=
\frac{1}{|\det T|}\int_{\RR^d}f(x)\,d\lambda_d(x).
\]
\end{proof}

\begin{corollary}[可積分関数に対する線形変数変換]\label{cor:linear-change-of-variables-integrable}
$T\in GL(d,\RR)$ とし，
$f:\RR^d\to\eRR$ を可積分関数とする．
このとき
\[
\int_{\RR^d}f(Tu)\,d\lambda_d(u)
=
\frac{1}{|\det T|}\int_{\RR^d}f(x)\,d\lambda_d(x)
\]
が成り立つ．
\end{corollary}
\begin{proof}
\cref{thm:linear-change-of-variables}を $f^+$ と $f^-$ に適用して差を取ればよい．
\end{proof}

\begin{corollary}[affine 変数変換]\label{cor:affine-change-of-variables}
$T\in GL(d,\RR)$，$b\in\RR^d$ とし，
$f:\RR^d\to[0,\infty]$ を非負可測関数とする．
このとき
\[
\int_{\RR^d}f(Tu+b)\,|\det T|\,d\lambda_d(u)
=
\int_{\RR^d}f(x)\,d\lambda_d(x)
\]
が成り立つ．
\end{corollary}
\begin{proof}
まず
\[
g(x):=f(x+b)
\]
とおくと，平行移動不変性の積分版
\[
\int_{\RR^d}g(x)\,d\lambda_d(x)=\int_{\RR^d}f(x)\,d\lambda_d(x)
\]
が成り立つ．
これは1変数の affine 変数変換（\cref{thm:one-dimensional-affine-change-variables}）を
各座標に順に適用し，\cref{thm:tonelli}を用いれば従う．

したがって線形変数変換（\cref{thm:linear-change-of-variables}）を $g$ に適用して
\[
\int_{\RR^d}f(Tu+b)\,d\lambda_d(u)
=
\int_{\RR^d}g(Tu)\,d\lambda_d(u)
=
\frac1{|\det T|}\int_{\RR^d}g(x)\,d\lambda_d(x)
=
\frac1{|\det T|}\int_{\RR^d}f(x)\,d\lambda_d(x)
\]
を得る．
両辺に $|\det T|$ を掛ければよい．
\end{proof}

\begin{corollary}[可積分関数に対する affine 変数変換]\label{cor:affine-change-of-variables-integrable}
$T\in GL(d,\RR)$，$b\in\RR^d$ とし，
$f:\RR^d\to\eRR$ を可積分関数とする．
このとき
\[
\int_{\RR^d}f(Tu+b)\,|\det T|\,d\lambda_d(u)
=
\int_{\RR^d}f(x)\,d\lambda_d(x)
\]
が成り立つ．
\end{corollary}
\begin{proof}
\cref{cor:affine-change-of-variables}を $f^+$ と $f^-$ に適用して差を取ればよい．
\end{proof}

\begin{corollary}[集合の体積の変換]\label{cor:volume-under-affine-map}
$T\in GL(d,\RR)$，$b\in\RR^d$ とし，
$A\subset\RR^d$ をLebesgue可測集合とする．
このとき
\[
\lambda_d(TA+b)=|\det T|\,\lambda_d(A)
\]
が成り立つ．
\end{corollary}
\begin{proof}
\cref{cor:affine-change-of-variables}を
\[
f=1_A
\]
に適用すれば
\[
\int_{\RR^d}1_A(Tu+b)\,|\det T|\,d\lambda_d(u)
=
\int_{\RR^d}1_A(x)\,d\lambda_d(x)
=
\lambda_d(A).
\]
一方，
\[
1_A(Tu+b)=1_{T^{-1}(A-b)}(u)
\]
だから左辺は
\[
|\det T|\,\lambda_d(T^{-1}(A-b)).
\]
ここで
\[
T^{-1}(A-b)
\]
を $B$ とおけば $A=TB+b$ であるから
\[
\lambda_d(A)=|\det T|\,\lambda_d(B).
\]
すなわち
\[
\lambda_d(TA+b)=|\det T|\,\lambda_d(A)
\]
である．
\end{proof}

\begin{example}[平行四辺形の面積]
\[
T=
\begin{pmatrix}
a & c\\
b & d
\end{pmatrix}
\]
とし，
\[
Q:=[0,1]^2
\]
とする．
このとき $TQ$ は原点を頂点にもつ平行四辺形であり，
その面積は
\[
\lambda_2(TQ)=|ad-bc|
\]
である．
\end{example}
\begin{proof}
\cref{cor:volume-under-affine-map}より
\[
\lambda_2(TQ)=|\det T|\,\lambda_2(Q)=|ad-bc|\cdot 1=|ad-bc|.
\]
\end{proof}

\section{解釈とまとめ}

\cref{thm:fubini}は，
\[
\iint f(x,y)\,dx\,dy
\]
を
\[
\int\left(\int f(x,y)\,dy\right)dx
\]
として計算してよい条件を与える定理であった．
非負関数なら\cref{thm:tonelli}，
符号付き関数なら
\[
\iint |f|<\infty
\]
という絶対可積分性の仮定のもとで\cref{thm:fubini}が使える．

また\cref{cor:affine-change-of-variables,cor:affine-change-of-variables-integrable}の変数変換公式は，
\[
\text{図形の面積・体積がどれだけ伸び縮みするか}
\]
を
\[
|\det T|
\]
で表す定理であった．
特に affine 変換
\[
x=Tu+b
\]
に対して
\[
dx=|\det T|\,du
\]
と見なしてよいことが，厳密に証明された．

これら2つの道具は，多変数積分を実際に計算する際の標準手順である．
今後は
\[
\text{積分順序を変える},\qquad \text{変数を取り替える}
\]
という操作を，単なる計算テクニックではなく，
Lebesgue積分論の定理として正当化して使えるようになる．
